
\chapter{Convergence of funcoids}


\section{Convergence}

Use Definition~\ref{def-osa-convergence} on the action of funcoids
on filters. In particular, $f$ converges to $\mathcal A$ regarding
$\mu$ on $\mathcal B$ iff
$\supfun{f}\mathcal B\sqsubseteq\supfun{\mu}\mathcal A$.
For a reloid $f$ and a funcoid $\mu$, apply this definition to
$\tofcd f$. This gives exactly the test using the reloid image,
since $\im(\tofcd f)=\im f$ by Section~\ref{sec-reloid-action}.

The convergence result under local continuity in
Section~\ref{sec-osa-convergence} uses the square restriction of
$\mu$ at $b=\supfun{\mu}\mathcal A$. For a funcoid this is
$\mu\sqcap(b\times^{\mathsf{FCD}}b)$ by
Proposition~\ref{prop-osa-rectangle-meet}. Thus, if
$g|_b$ is continuous from this square restriction to $\nu$
and $\mathcal A\sqsubseteq b$, then postcomposition by $g$
takes convergence to $\mathcal A$ regarding $\mu$ to convergence
to $\supfun{g}\mathcal A$ regarding $\nu$.
For unrestricted continuity, Proposition~\ref{prop-osa-cont-conv}
gives the result without the inequality $\mathcal A\sqsubseteq b$.

\section{Relationships between convergence and continuity}
For a funcoid $f$ from $\Ob\mu$ to $\Ob\nu$, the local
continuity condition $f\in\continuous(\mu|_{\mathcal A},\nu)$
implies that $f|_{\supfun{\mu}\mathcal A}$ converges regarding
$\nu$ to $\supfun{f}\mathcal A$. If $\mathcal A$ is an
ultrafilter, the converse holds as well. Indeed, ultrafilters
are atoms of the filter lattice and the funcoid action is
ordered by elements and preserves the least element, so the
local criterion of Section~\ref{sec-osa-convergence} applies.
Unrestricted continuity implies the same convergence by
Proposition~\ref{prop-osa-cont-conv}.

\section{Convergence of joins}
The join criteria of Section~\ref{sec-osa-convergence} apply here.
The action preserves binary joins of filters, and images preserve
all joins of parallel funcoids by Corollary~\ref{fcd-dom-join}.

\section{Limit}
Take the points in Definition~\ref{def-osa-limit} to be the principal
singleton filters. We write the value of $\lim^\mu f$ as the
underlying point $a$ when the limit filter is $\uparrow^{\Src\mu}\{a\}$.
If $f$ is a non-least funcoid, its image is non-least, so the
uniqueness result in Section~\ref{sec-osa-separation} applies when
$\mu$ is $T_2$-separable. The notation $\lim_{\mathcal B}^\mu f$
then uses the limit of the restriction to the filter $\mathcal B$.
For a reloid, the action through $\tofcd$ provides the same notion.

\section{Generalized limit}

\subsection{\index{limit!generalized}Definition}

Let $\mu$ and $\nu$ be endofuncoids. Let $G$ be a transitive permutation
group on $\Ob\mu$.

For an element $r\in G$ we will denote $\uparrow r=\uparrow^{\mathsf{FCD}(\Ob\mu,\Ob\mu)}r$.

We require that $\mu$ and every $r\in G$ commute, that is
\[
\mu\circ\uparrow r=\uparrow r\circ\mu.
\]


We require for every $y\in\Ob\nu$ 
\begin{equation}
\nu\sqsupseteq\rsupfun{\nu}@\{y\}\times^{\mathsf{FCD}}\rsupfun{\nu}@\{y\}.\label{lim-squares}
\end{equation}

\begin{prop}
Formula (\ref{lim-squares}) follows from $\nu\sqsupseteq\nu\circ\nu^{-1}$.\end{prop}
\begin{proof}
Let $\nu\sqsupseteq\nu\circ\nu^{-1}$. Then
\begin{align*}
\rsupfun{\nu}@\{y\}\times^{\mathsf{FCD}}\rsupfun{\nu}@\{y\} & =\\
\supfun{\nu}@\{y\}\times^{\mathsf{FCD}}\supfun{\nu}@\{y\} & =\\
\nu\circ(\uparrow^{\Ob\nu}\{y\}\times^{\mathsf{FCD}}\uparrow^{\Ob\nu}\{y\})\circ\nu^{-1} & =\\
\nu\circ\uparrow^{\mathsf{FCD}(\Ob\nu,\Ob\nu)}(\{y\}\times\{y\})\circ\nu^{-1} & \sqsubseteq\\
\nu\circ1_{\Ob\nu}^{\mathsf{FCD}}\circ\nu^{-1} & =\\
\nu\circ\nu^{-1} & \sqsubseteq\nu.
\end{align*}
\end{proof}
\begin{rem}
The formula (\ref{lim-squares}) usually works if $\nu$ is a proximity.
It does not work if $\mu$ is a pretopology or preclosure.
\end{rem}
We are going to consider (generalized) limits of arbitrary functions
acting from $\Ob\mu$ to $\Ob\nu$. (The functions in consideration
are not required to be continuous.)
\begin{rem}
Most typically $G$ is the group of translations of some topological
vector space.
\end{rem}
Apply Definition~\ref{def-osa-xlim} with
$\rho(r)=\uparrow r$. This is a group homomorphism because the
principal funcoid construction is a functor. At a point $x$ of the
underlying set $\Ob\mu$, use the principal singleton filter
$\uparrow^{\Ob\mu}\{x\}$ as the action element in
Definition~\ref{def-osa-xlim}. Thus $\xlim_x f$ is
$\xlim(f|_{\rsupfun{\mu}@\{x\}})$. We assume
$\dom f\sqsupseteq\rsupfun{\mu}@\{x\}$ for the results below.
Here $\xlim_x f$ is unchanged on replacing $\mu$ by $\Compl\mu$
because their action values on principal singleton filters agree.

The function $\tau$ will define an injection from the set of points
of the space $\nu$ (``numbers'', ``points'', or ``vectors'')
to the set of all (generalized) limits (i.e. values which $\xlim_{x}f$
may take).
\begin{defn}
For $y\in\Ob\nu$, put
$\tau(y)\eqdef\setcond{\rsupfun{\mu}@\{x\}\times^{\mathsf{FCD}}\rsupfun{\nu}@\{y\}}{x\in\Ob\mu}$.\end{defn}
\begin{prop}
$\tau(y)=\setcond{(\rsupfun{\mu}@\{x\}\times^{\mathsf{FCD}}\rsupfun{\nu}@\{y\})\circ\uparrow r}{r\in G}$
for every fixed $x\in\Ob\mu$.\end{prop}
\begin{proof}
~
\begin{align*}
(\rsupfun{\mu}@\{x\}\times^{\mathsf{FCD}}\rsupfun{\nu}@\{y\})\circ\uparrow r & =\\
\supfun{\uparrow r^{-1}}\rsupfun{\mu}@\{x\}\times^{\mathsf{FCD}}\rsupfun{\nu}@\{y\} & =\\
\supfun{\mu}\rsupfun{\uparrow r^{-1}}@\{x\}\times^{\mathsf{FCD}}\rsupfun{\nu}@\{y\} & =\\
\rsupfun{\mu}@\{r^{-1}x\}\times^{\mathsf{FCD}}\rsupfun{\nu}@\{y\} & \in \setcond{\rsupfun{\mu}@\{x\}\times^{\mathsf{FCD}}\rsupfun{\nu}@\{y\}}{x\in\Ob\mu}.
\end{align*}


Reversely $\rsupfun{\mu}@\{x\}\times^{\mathsf{FCD}}\rsupfun{\nu}@\{y\}=(\rsupfun{\mu}@\{x\}\times^{\mathsf{FCD}}\rsupfun{\nu}@\{y\})\circ\uparrow e$
where $e$ is the identify element of $G$.\end{proof}
\begin{prop}
$\tau(y)=\xlim(\rsupfun{\mu}@\{x\}\times^{\mathsf{FCD}}\uparrow^{\Base(\Ob\nu)}\{y\})$
(for every $x$). Informally: Every $\tau(y)$ is a generalized limit
of a constant funcoid.\end{prop}
\begin{proof}
~
\begin{align*}
\xlim(\rsupfun{\mu}@\{x\}\times^{\mathsf{FCD}}\uparrow^{\Base(\Ob\nu)}\{y\}) & =\\
\setcond{\nu\circ(\rsupfun{\mu}@\{x\}\times^{\mathsf{FCD}}\uparrow^{\Base(\Ob\nu)}\{y\})\circ\uparrow r}{r\in G} & =\\
\setcond{(\rsupfun{\mu}@\{x\}\times^{\mathsf{FCD}}\rsupfun{\nu}@\{y\})\circ\uparrow r}{r\in G} & =\tau(y).
\end{align*}
\end{proof}
\begin{thm}
If $f$ is a function and $f|_{\rsupfun{\mu}@\{x\}}\in\continuous(\mu,\nu)$
and $\rsupfun{\mu}@\{x\}\sqsupseteq\uparrow^{\Ob\mu}\{x\}$
then $\xlim_{x}f=\tau(fx)$.\end{thm}
\begin{proof}
$f|_{\rsupfun{\mu}@\{x\}}\circ\mu\sqsubseteq\nu\circ f|_{\rsupfun{\mu}@\{x\}}\sqsubseteq\nu\circ f$;
thus $\langle f\rangle\rsupfun{\mu}@\{x\}\sqsubseteq\supfun{\nu}\rsupfun{f}@\{x\}$;
consequently we have
\begin{gather*}
\nu\sqsupseteq\supfun{\nu}\rsupfun{f}@\{x\}\times^{\mathsf{FCD}}\supfun{\nu}\rsupfun{f}@\{x\}\sqsupseteq\supfun f\rsupfun{\mu}@\{x\}\times^{\mathsf{FCD}}\supfun{\nu}\rsupfun{f}@\{x\}.\\
\begin{aligned}\nu\circ f|_{\rsupfun{\mu}@\{x\}} & \sqsupseteq\\
(\supfun f\rsupfun{\mu}@\{x\}\times^{\mathsf{FCD}}\supfun{\nu}\rsupfun{f}@\{x\})\circ f|_{\rsupfun{\mu}@\{x\}} & =\\
(f|_{\rsupfun{\mu}\{x\}})^{-1}\supfun f\rsupfun{\mu}@\{x\}\times^{\mathsf{FCD}}\supfun{\nu}\rsupfun{f}@\{x\} & \sqsupseteq\\
\supfun{\id_{\dom f|_{\rsupfun{\mu}\{x\}}}^{\mathsf{FCD}}}\rsupfun{\mu}@\{x\}\times^{\mathsf{FCD}}\supfun{\nu}\rsupfun{f}@\{x\} & \sqsupseteq\\
\dom f|_{\rsupfun{\mu}@\{x\}}\times^{\mathsf{FCD}}\supfun{\nu}\rsupfun{f}@\{x\} & =\\
\rsupfun{\mu}@\{x\}\times^{\mathsf{FCD}}\supfun{\nu}\rsupfun{f}@\{x\}.
\end{aligned}
\end{gather*}


$\im(\nu\circ f|_{\rsupfun{\mu}@\{x\}})=\supfun{\nu}\rsupfun{f}@\{x\}$;
\begin{align*}
\nu\circ f|_{\rsupfun{\mu}@\{x\}} & \sqsubseteq\\
\rsupfun{\mu}@\{x\}\times^{\mathsf{FCD}}\im(\nu\circ f|_{\rsupfun{\mu}@\{x\}}) & =\\
\rsupfun{\mu}@\{x\}\times^{\mathsf{FCD}}\supfun{\nu}\rsupfun{f}@\{x\}.
\end{align*}
So $\nu\circ f|_{\rsupfun{\mu}@\{x\}}=\rsupfun{\mu}@\{x\}\times^{\mathsf{FCD}}\supfun{\nu}\rsupfun{f}@\{x\}$.

Thus $\xlim_{x}f=\setcond{(\rsupfun{\mu}@\{x\}\times^{\mathsf{FCD}}\supfun{\nu}\rsupfun{f}@\{x\})\circ\uparrow r}{r\in G}=\tau(fx)$.\end{proof}
\begin{rem}
Without the requirement of $\rsupfun{\mu}@\{x\}\sqsupseteq\uparrow^{\Ob\mu}\{x\}$
the last theorem would not work in the case of removable singularity.\end{rem}
\begin{thm}
Let $\nu\sqsubseteq\nu\circ\nu$. If $f|_{\rsupfun{\mu}@\{x\}}\overset{\nu}{\rightarrow}\uparrow^{\Ob\mu}\{y\}$
then $\xlim_{x}f=\tau(y)$.\end{thm}
\begin{proof}
$\im f|_{\rsupfun{\mu}@\{x\}}\sqsubseteq\rsupfun{\nu}@\{y\}$;
$\supfun f\rsupfun{\mu}@\{x\}\sqsubseteq\rsupfun{\nu}@\{y\}$;
\begin{align*}
\nu\circ f|_{\rsupfun{\mu}@\{x\}} & \sqsupseteq\\
(\rsupfun{\nu}@\{y\}\times^{\mathsf{FCD}}\rsupfun{\nu}@\{y\})\circ f|_{\rsupfun{\mu}@\{x\}} & =\\
\supfun{(f|_{\rsupfun{\mu}@\{x\}})^{-1}}\rsupfun{\nu}@\{y\}\times^{\mathsf{FCD}}\rsupfun{\nu}@\{y\} & =\\
\supfun{\id_{\rsupfun{\mu}\{x\}}^{\mathsf{FCD}}\circ f^{-1}}\rsupfun{\nu}@\{y\}\times^{\mathsf{FCD}}\rsupfun{\nu}@\{y\} & \sqsupseteq\\
\supfun{\id_{\rsupfun{\mu}\{x\}}^{\mathsf{FCD}}\circ f^{-1}}\supfun f\rsupfun{\mu}@\{x\}\times^{\mathsf{FCD}}\rsupfun{\nu}@\{y\} & =\\
\supfun{\id_{\rsupfun{\mu}\{x\}}^{\mathsf{FCD}}}\supfun{f^{-1}\circ f}\rsupfun{\mu}@\{x\}\times^{\mathsf{FCD}}\rsupfun{\nu}@\{y\} & \sqsupseteq\\
\supfun{\id_{\rsupfun{\mu}\{x\}}^{\mathsf{FCD}}}\supfun{\id_{\rsupfun{\mu}\{x\}}^{\mathsf{FCD}}}\rsupfun{\mu}@\{x\}\times^{\mathsf{FCD}}\rsupfun{\nu}@\{y\} & =\\
\rsupfun{\mu}@\{x\}\times^{\mathsf{FCD}}\rsupfun{\nu}@\{y\}.
\end{align*}


On the other hand, $f|_{\rsupfun{\mu}@\{x\}}\sqsubseteq\rsupfun{\mu}@\{x\}\times^{\mathsf{FCD}}\rsupfun{\nu}@\{y\}$;

$\nu\circ f|_{\rsupfun{\mu}@\{x\}}\sqsubseteq\rsupfun{\mu}@\{x\}\times^{\mathsf{FCD}}\supfun{\nu}\rsupfun{\nu}@\{y\}\sqsubseteq\rsupfun{\mu}@\{x\}\times^{\mathsf{FCD}}\rsupfun{\nu}@\{y\}$.

So $\nu\circ f|_{\rsupfun{\mu}@\{x\}}=\rsupfun{\mu}@\{x\}\times^{\mathsf{FCD}}\rsupfun{\nu}@\{y\}$.

$\xlim_{x}f=\setcond{\nu\circ f|_{\rsupfun{\mu}@\{x\}}\circ\uparrow r}{r\in G}=\setcond{(\rsupfun{\mu}@\{x\}\times^{\mathsf{FCD}}\rsupfun{\nu}@\{y\})\circ\uparrow r}{r\in G}=\tau(y)$.\end{proof}
\begin{cor}
If $\lim_{\rsupfun{\mu}@\{x\}}^{\nu}f=y$ then $\xlim_{x}f=\tau(y)$ (provided that $\nu\sqsubseteq\nu\circ\nu$).
\end{cor}
We have injective $\tau$ if $\rsupfun{\nu}@\{y_{1}\}\sqcap\rsupfun{\nu}@\{y_{2}\}=\bot^{\mathscr{F}(\Ob\mu)}$
for every distinct $y_{1},y_{2}\in\Ob\nu$ that is if $\nu$ is $T_{2}$-separable.

\section{Expressing limits as implications}

When you studied limits in the school, you was told that
$\lim_{x\to\alpha} f(x) = \beta$ when $x\to\alpha$ implies
$f(x) \to \beta$. Now let us formalize this.

\begin{prop}
  The following are pairwise equivalent for funcoids $\mu$, $\nu$, $f$ of
  suitable (``compatible'') sources and destinations:
  \begin{enumerate}
    \item\label{implim-one} $f|_{\langle \mu \rangle^{\ast} \{ \alpha \}}
    \overset{\nu}{\rightarrow} \beta$;
    
    \item\label{implim-f} $\forall x \in \mathscr{F} (\Ob \mu) : \left( x
    \overset{\mu}{\rightarrow} \alpha \Rightarrow \supfun{f} x
    \overset{\nu}{\rightarrow} \beta \right)$;
    
    \item\label{implim-a} $\forall x \in \atoms^{\mathscr{F} (\Ob \mu)} : \left( x
    \overset{\mu}{\rightarrow} \alpha \Rightarrow \supfun{f} x
    \overset{\nu}{\rightarrow} \beta \right)$.
  \end{enumerate}
\end{prop}

\begin{proof}
  ~
  \begin{description}
    \item[\ref{implim-one}$\Leftrightarrow$\ref{implim-f}]
    \begin{multline*}
    \forall x \in \mathscr{F} (\Ob \mu) :
    \left( x \overset{\mu}{\rightarrow} \alpha \Rightarrow \supfun{f} x
    \overset{\nu}{\rightarrow} \beta \right) \Leftrightarrow \\ \forall x \in
    \mathscr{F} (\Ob \mu) : (x \sqsubseteq \langle \mu \rangle \alpha
    \Rightarrow \\ \supfun{f} x \sqsubseteq \supfun{\nu} \beta)
    \Leftrightarrow \supfun{f} \langle \mu \rangle \alpha \sqsubseteq
    \supfun{\nu} \beta \Leftrightarrow f|_{\langle \mu \rangle^{\ast}
    \{ \alpha \}} \overset{\nu}{\rightarrow} \beta.
    \end{multline*}
    
    \item[\ref{implim-f}$\Rightarrow$\ref{implim-a}] Obvious.
    
    \item[\ref{implim-a}$\Rightarrow$\ref{implim-f}] Let \ref{implim-a} hold. Then for $x \in \mathscr{F}
    (\Ob \mu)$ we have $x \overset{\mu}{\rightarrow} \alpha
    \Leftrightarrow x \sqsubseteq \langle \mu \rangle \alpha \Leftrightarrow
    \forall x' \in \atoms x : x' \sqsubseteq \langle \mu \rangle \alpha
    \Leftrightarrow \forall x' \in \atoms x : x'
    \overset{\mu}{\rightarrow} \alpha \Rightarrow \forall x' \in \atoms
    x : \supfun{f} x' \overset{\nu}{\rightarrow} \beta \Leftrightarrow
    \forall x' \in \atoms x : \left\langle f \right\rangle x'
    \sqsubseteq \supfun{\nu} \beta \Leftrightarrow \bigsqcup_{x' \in
    \atoms x} \supfun{f} x' \sqsubseteq \supfun{\nu} \beta
    \Leftrightarrow \supfun{f} x \sqsubseteq \supfun{\nu} \beta
    \Leftrightarrow \supfun{f} x \overset{\nu}{\rightarrow} \beta$.
  \end{description}
\end{proof}

\begin{lem}
  If $f$ is an enterely defined monovalued funcoid and $x$ is an ultrafilter,
  $y$ is a filter, then $\supfun{f} x \sqsubseteq y \Leftrightarrow x
  \sqsubseteq \langle f^{- 1} \rangle y$.
\end{lem}

\begin{proof}
  $\supfun{f} x$ is an ultrafilter. $\supfun{f} x \sqsubseteq y
  \Leftrightarrow \supfun{f} x \nasymp y \Leftrightarrow x \nasymp
  \langle f^{- 1} \rangle y \Leftrightarrow x \sqsubseteq \langle f^{- 1}
  \rangle y$.
\end{proof}

\begin{prop}
  The following are pairwise equivalent for funcoids $\mu$, $\nu$, $f$, $g$ of
  suitable (``compatible'') sources and destinations provided that $g$ is
  entirely defined and monovalued:
  \begin{enumerate}
    \item\label{limback-one} $(f \circ g^{- 1}) |_{\langle \mu \rangle^{\ast} \{ \alpha \}}
    \overset{\nu}{\rightarrow} \beta$;
    
    \item\label{limback-f} $\forall x \in \mathscr{F} (\Ob \mu) : \left( \langle g
    \rangle x \overset{\mu}{\rightarrow} \alpha \Rightarrow \supfun{f}
    x \overset{\nu}{\rightarrow} \beta \right)$;
    
    \item\label{limback-a} $\forall x \in \atoms^{\mathscr{F} (\Ob \mu)} : \left(
    \supfun{g} x \overset{\mu}{\rightarrow} \alpha \Rightarrow \langle
    f \rangle x \overset{\nu}{\rightarrow} \beta \right)$.
  \end{enumerate}
\end{prop}

\begin{proof}
  ~
  \begin{description}
    \item[\ref{limback-one}$\Leftrightarrow$\ref{limback-a}] Equivalently transforming: $(f \circ g^{- 1})
    |_{\langle \mu \rangle^{\ast} \{ \alpha \}} \overset{\nu}{\rightarrow}
    \beta$; $\supfun{f} \langle g^{- 1} \rangle^{\ast} \langle \mu
    \rangle^{\ast} \{ \alpha \} \sqsubseteq \supfun{\nu}^{\ast} \{
    \beta \}$; for every $x \in \atoms^{\mathscr{F} (\Ob \mu)}$ we
    have $x \sqsubseteq \langle g^{- 1} \rangle \langle \mu \rangle^{\ast} \{
    \alpha \} \Rightarrow \supfun{f} x \sqsubseteq \langle \nu
    \rangle^{\ast} \{ \beta \}$; what by the lemma is equivalent to $\langle g
    \rangle x \sqsubseteq \langle \mu \rangle^{\ast} \{ \alpha \} \Rightarrow
    \supfun{f} x \sqsubseteq \supfun{\nu}^{\ast} \{ \beta \}$
    that is $\supfun{g} x \overset{\mu}{\rightarrow} \alpha \Rightarrow
    \supfun{f} x \overset{\nu}{\rightarrow} \beta$.
    
    \item[\ref{limback-a}$\Leftrightarrow$\ref{limback-f}] Let $x \in \mathscr{F} (\Ob \mu)$ and~\ref{limback-a}
    holds. Let $\supfun{g} x \overset{\mu}{\rightarrow} \alpha$. Then
    $\forall x' \in \atoms x : \supfun{g} x'
    \overset{\mu}{\rightarrow} \alpha$ and thus $\supfun{f} x'
    \overset{\nu}{\rightarrow} \beta$ that is $\supfun{f} x'
    \sqsubseteq \supfun{\nu} \beta$. $\supfun{f} x =
    \bigsqcup_{x' \in \atoms x} \supfun{f} x' \sqsubseteq \langle
    \nu \rangle \beta$ that is $\supfun{f} x \overset{\nu}{\rightarrow}
    \beta$.
  \end{description}
\end{proof}

\begin{problem}
  Can the theorem be strenhtened for: a. non-monovalued; b. not entirely defined $g$?
  (The problem seems easy but I have not checked it.)
\end{problem}
